\documentclass[conference,10pt]{IEEEtran}
\usepackage{cite}
\usepackage{amsmath,amssymb,amsfonts}
\usepackage{algorithmic}
\usepackage{graphicx}
\usepackage{textcomp}
\usepackage{xcolor}
\usepackage{booktabs}
\usepackage{multirow}
\usepackage{float}
\usepackage{hyperref}
\usepackage[turkish]{babel}
\usepackage[utf8]{inputenc}

\begin{document}

\title{Uydu Telemetri Sistemleri için Uçtan Uca Makine Öğrenmesi Tabanlı Anomali Tespiti, Canlı İzleme ve Karar Destek Ağı: Kapsamlı Bitirme Projesi Raporu\\}

\author{\IEEEauthorblockN{Proje Ekibi / Öğrenci Adı}
\IEEEauthorblockA{\textit{Mühendislik Fakültesi} \\
\textit{Üniversite/Kurum Adı}\\
Şehir, Ülke \\
ogrenci@domain.com}
}

\maketitle

\begin{abstract}
Uzay sistemlerindeki otonomi ihtiyacı, uydu alt sistemlerinden gelen devasa telemetri verilerinin gerçek zamanlı analizini zorunlu kılmaktadır. Geleneksel eşik tabanlı (Out-Of-Limits) arıza tespit sistemleri, sensörler arası gizli korelasyonları ve dinamik uzay koşullarını modellemede yetersiz kalmaktadır. Bu bitirme projesi kapsamında, ESA OPS-SAT uydusunun yönelim belirleme ve tepki tekerlekleri (reaction wheels) verilerinden beslenen, uçtan uca eksiksiz bir anomali tespit (FDIR) boru hattı ve etkileşimli bir web paneli (Dashboard) geliştirilmiştir. 

Proje; Veri İnceleme (EDA), Ön İşleme, Gelişmiş Özellik Mühendisliği (Zaman, Frekans, Fiziksel uzaylar), 12 farklı Yapay Zeka Modelinin Eğitimi (Gözetimli ve Gözetimsiz), Ablasyon Analizi, SHAP ile Açıklanabilir Yapay Zeka (XAI) ve Gerçek Zamanlı (Canlı) İzleme modüllerinin tamamını barındırmaktadır. Ayrıca, otomatik deney yürütme betikleri (create\_ablation\_nb.py, patch\_notebook.py), birim testler (pytest) ve model/veri metrik raporlama sistemleri (adv\_metrics.json) gibi yazılım mühendisliği pratikleri entegre edilmiştir.

Oluşturulan altyapıda, \textit{TelemetriPreprocessor} ve \textit{ReactionWheelFeatureEngineer} modülleri kullanılarak yüzlerce yeni istatistiksel ve fiziksel öznitelik üretilmiş, Varyans ve Korelasyon filtreleriyle optimal özellik seti elde edilmiştir. Gözetimli modellerde \textbf{MLP} \%99.2 AUC-ROC ve \%91.9 F1 skoruyla en yüksek performansı verirken; gözetimsiz (etiketsiz) veride derin öğrenme tabanlı \textbf{Autoencoder} \%89.3 AUC-ROC başarımı sağlamıştır. Geliştirilen 9 sekme/sayfadan oluşan Plotly Dash arayüzü sayesinde, kullanıcılar canlı telemetri akışını farklı hızlarda (1x, 5x, 20x) izleyebilmekte, anomali skorlarını dinamik olarak görebilmekte ve \texttt{reportlab} ile otomatik PDF raporu üretebilmektedir.
\end{abstract}

\begin{IEEEkeywords}
Uydu Telemetrisi, Makine Öğrenmesi, Anomali Tespiti, Otokodlayıcı, MLP, Dash, SHAP, Ablasyon, Yazılım Mühendisliği, ESA OPS-SAT.
\end{IEEEkeywords}

\section{Giriş ve Motivasyon}
Modern uydular, Yönelim Belirleme ve Kontrol Sistemi (ADCS), Termal Kontrol ve Güç Sistemlerinden saniyede binlerce veri noktası (telemetri) üretmektedir. Tepki tekerlekleri (Reaction Wheels) ve manyetometreler gibi kritik parçalardaki mekanik yorgunluk veya donanımsal anormallikler uydunun tüm görevini tehlikeye atabilir. Endüstri standardı olan OOL (Out-of-Limits) sistemleri yüksek oranlarda yanlış alarm üretmekte ve sensörler arası (multivariate) anomalileri gözden kaçırmaktadır.

Bu projenin temel vizyonu, literatürdeki teorik makine öğrenmesi modellerini test etmekle kalmayıp, veri yüklemesinden canlı izlemeye, otomatik Jupyter Notebook düzenlemelerinden PDF raporlamasına ve otomatik test aşamalarına kadar tamamen \textbf{kullanıma hazır, entegre bir ürün (Karar Destek Sistemi)} geliştirmektir. 

\section{Sistem Mimarisi ve Proje Dizini}
Proje, modüler bir MVC (Model-View-Controller) benzeri yapıda tasarlanmıştır:
\begin{itemize}
    \item \textbf{\texttt{src/} (Backend ve Veri Boru Hattı):} \texttt{data\_loader.py}, \texttt{preprocessor.py}, \texttt{feature\_engineer.py} ve \texttt{utils.py} ile veri dönüşümleri.
    \item \textbf{\texttt{app/} (Frontend ve Web Sunucusu):} Plotly Dash tabanlı \texttt{app.py}, \texttt{ablation\_page.py}, \texttt{assets/} (CSS stilleri) ve \texttt{utils/model\_loader.py}.
    \item \textbf{\texttt{notebooks/} (Ar-Ge ve Deneyler):} 01'den 08'e kadar sıralı veri inceleme, özellik üretimi, model eğitimi ve SHAP analizi dosyaları. Ek olarak \texttt{shap\_code.txt} ve catboost eğitim logları.
    \item \textbf{\texttt{models/} ve \texttt{reports/}:} Eğitilmiş modellerin (h5, keras, joblib) gözetimli ve gözetimsiz (\texttt{models/unsupervised/}) olarak ayrıştırılması. Metriklerin \texttt{final\_comparison.json} ve \texttt{adv\_metrics.json} ile saklanması.
    \item \textbf{\texttt{tests/} (Birim Testleri):} Yazılım kalitesini garantilemek için yazılan \texttt{test\_data\_loader.py} ve \texttt{test\_model\_loader.py} (\texttt{pytest} tabanlı).
    \item \textbf{Otomasyon Betikleri:} Sürekli entegrasyon (CI/CD) mantığıyla çalışan \texttt{create\_ablation\_nb.py}, \texttt{patch\_notebook.py} ve \texttt{revert\_notebook.py} ile notebook'ların programatik olarak oluşturulması ve düzenlenmesi.
\end{itemize}

\section{Veri Yükleme ve Gelişmiş Ön İşleme}
\subsection{Evrensel Veri Yükleyici (DataLoader)}
\texttt{src/data\_loader.py} modülü; \textbf{CSV, Parquet, JSON, H5, HDF5 ve Excel} formatlarındaki tüm telemetri dosyalarını dinamik okuyabilmektedir. \texttt{validate\_data} ve \texttt{get\_summary} metodları ile yüklenen verinin eksik kolon kontrolü ve bellek analizi (MB bazında) anlık yapılmaktadır.

\subsection{Telemetri Ön İşleme Motoru (Preprocessor)}
Uydu sensörlerindeki termal dalgalanmalar, kozmik radyasyon sıçramaları ve bağlantı kopmaları (packet loss) \texttt{src/preprocessor.py} ile bertaraf edilmiştir:
\begin{itemize}
    \item \textbf{Eksik Veri Doldurma:} İleri/Geri Tamamlama (ffill/bfill), Doğrusal İnterpolasyon, 3. Derece Kübik Spline İnterpolasyonu ve \textit{k}-En Yakın Komşu ($k=5$ KNN) yöntemleri algoritmik olarak seçilebilmektedir.
    \item \textbf{Sinyal Filtreleme:} Düşük frekanslı gerçek hareketleri koruyup yüksek frekanslı gürültüyü atmak için 3. derece \textbf{Savitzky-Golay} filtresi, 4. derece \textbf{Butterworth Low-Pass} filtresi (0.2 cutoff) ve \textbf{Medyan (Median)} filtreleri uygulanmıştır.
    \item \textbf{Aykırı Değer (Outlier) Kırpma:} Çeyrekler Arası Aralık (IQR), Z-Skoru ve sağlam (robust) \textbf{Modifiye Z-Skoru} (MAD tabanlı) kullanılarak anormal uç değerler, dinamik alt/üst sınırlarla kırpılmıştır (clipping).
    \item \textbf{Ölçeklendirme:} MinMaxScaler, StandardScaler ve RobustScaler. Eğitilen scaler dosyaları \texttt{save\_scaler} metodu ile \texttt{joblib} kullanılarak kaydedilir.
\end{itemize}

\section{Kapsamlı Özellik Mühendisliği (Feature Engineering)}
Saf sensör okumaları tek başına arıza örüntüsü barındırmayabilir. \texttt{src/feature\_engineer.py} ve \texttt{app/utils/feature\_extractor.py} ile yüzlerce sentetik özellik türetilmiştir:

\subsection{Zaman Alanı (Time Domain) ve Enerji Metrikleri}
Tekerleklerin son durumunu anlamak için 30, 60 ve 120 saniyelik yuvarlanan pencereler (rolling windows) üzerinden ortalama, varyans, standart sapma, basıklık (kurtosis), çarpıklık (skewness) ve IQR hesaplanmıştır. Enerji metrikleri şunlardır:
\begin{itemize}
    \item \textbf{Kök Ortalama Kare (RMS):} Sinyal gücü ($RMS = \sqrt{\frac{1}{w}\sum x^2}$).
    \item \textbf{Tepe Faktörü (Crest Factor):} Anlık darbelerin göstergesi ($CF = \frac{x_{max}}{RMS}$).
    \item \textbf{Peak-to-Peak (P2P):} Pencere içi genlik salınım aralığı.
\end{itemize}

\subsection{Frekans ve Türev (Kinetik) Özellikler}
Dinamik hareketlerin tespiti için 1. türev olan Değişim Oranı (\textbf{Rate of Change - RoC}) ve 2. türev olan Sarsıntı (\textbf{Jerk}) diferansiyel denklemlerle hesaplanmıştır. Titreşim ölçümü için \textbf{Sıfır-Geçiş Oranı (ZCR)} üretilmiştir. Zamansal Gecikme (Lag) özellikleri ile (1, 5, 10, 30 ve 60 adım öncesi) RNN mimarilerine ihtiyaç duyulmadan tarihsel bağlam sisteme eklenmiştir.

\subsection{Fiziksel / Çok Değişkenli Özellikler (Domain Özgü)}
\begin{itemize}
    \item \textbf{Sensör Sentezi:} ESA OPS-SAT manyetometrelerinin (CADC0872, 0873, 0874) vektör şiddeti (MAG\_Magnitude = $\sqrt{x^2+y^2+z^2}$) ve uyumsuzluk farkı (MAG\_Sync\_Diff) formüle edilmiştir. Fotodiyotlar (PD) için toplam ışıma (PD\_Total\_Sum) ve varyans (PD\_Max\_Variance) bulunmuştur.
    \item \textbf{PCA ve Mahalanobis:} Çok boyutlu uzayda varyansı açıklayan 5 Temel Bileşen (PCA) oluşturulmuştur. Tüm sistemi tek skora indiren \textbf{Mahalanobis Uzaklığı}, kovaryans matrisinin sözde-tersi (pseudo-inverse) üzerinden çözülerek sisteme eklenmiştir.
    \item \textbf{Segment Bazlı Özellik Çıkarımı:} Ham veriler `extract_segment_features` metodu ile zaman serisi segmentleri etrafında gruplanarak Custom RMS, Custom P2P ve Custom ZCR metriklerine indirgenmiştir.
\end{itemize}
Boyutluluk lanetini önlemek için \textbf{Variance Threshold} ($1e^{-5}$) ve \textbf{Correlation Threshold} (Pearson $>0.95$) algoritmaları gereksiz değişkenleri silmiştir.

\section{Modelleme ve Algoritma Çeşitliliği}
\begin{enumerate}
    \item \textbf{Gözetimli Modeller:} Etiketli arıza verisiyle sınıf dengesizliklerine (\texttt{imbalanced-learn}) karşı korumalı \textbf{Random Forest, XGBoost, CatBoost, LightGBM, SVM} ve Keras altyapılı \textbf{MLP} eğitilmiştir. Modellerin ortak karar alması için bir \textbf{Stacking Ensemble} oluşturulmuştur.
    \item \textbf{Gözetimsiz Modeller:} \textbf{Isolation Forest}, \textbf{OneClassSVM}, \textbf{LOF}, \textbf{K-Means Clustering} ve derin öğrenme \textbf{Autoencoder} ile etiket olmadan yeniden yapılandırma (reconstruction) veya sınır-uzaklık hesaplarıyla aykırı değer tespiti yapılmıştır.
\end{enumerate}

\section{Gerçek Zamanlı Karar Destek Arayüzü (Web Dashboard)}
\texttt{app/app.py} web arayüzü tam 9 ana sayfadan/işlevden oluşur:

\subsection{1. Dashboard ve İstatistikler}
Sistemin genel sağlık matrisini, ısı haritası (Heatmap) üzerinden Accuracy, Precision, Recall, AUC skorları kıyaslamasını, F1 Skoru sıralama grafiklerini ve sistem operasyon loglarını sunar.
\subsection{2. Veri Yükleme ve Yabancı Veri Uyumu}
Kullanıcıların kendi verilerini sürükle-bırak ile yükleyebildiği alandır. \textbf{Auto-Adaptation} mekanizması: Eğer yüklenen CSV'de gerekli özellikler eksikse, \texttt{app/utils/feature\_extractor.py} tetiklenir, ham sensör değerlerinden otomatik olarak RMS, PCA ve Lag özellikleri çalışma zamanında üretilip analize hazır hale getirilir. Zaman serisi grafikleri çizdirilir.
\subsection{3. Canlı İzleme (Live Stream Simulation)}
"Canlı İzleme" sayfası, uzaydan anlık veri inişini simüle eder. Hız (1x, 5x, 20x) ve model kontrolleri içerir. Eşik aşıldığında "Telemetri Sinyali" kırmızı anomali noktaları oluşturur ve sistem \textbf{Alarm Paneli}ne anında log düşer.
\subsection{4. Analiz}
Gözetimli ve gözetimsiz modellerin çoklu seçimle aynı anda koşturulmasını (multi-model processing) sağlar. "Eşik Çarpanı" (Threshold Multiplier 0.5 ile 1.5 arası) kaydırıcısı ile algoritma hassasiyetine manuel müdahale izni verilir. İlerleme çubuklarıyla (Progress bar) sonuç üretir.
\subsection{5. Sonuçlar (Ensemble Çıktısı)}
Birden fazla modelin skorları kombine edilerek ikili (Binary Ensemble) ve sürekli (Continuous Severity Score) bir birleşim skoru üretilir. Hangi segmentin arızalı olduğu liste halinde verilir.
\subsection{6. Model Performansı}
Plotly kütüphanesiyle dinamik olarak \textbf{ROC Eğrileri} (ROC Curves) ve tüm metriklerin çok yönlü mukayesesini sağlayan \textbf{Radar (Spider) Çizelgesi} oluşturur. Dinamik renklendirilen "Metrik Tablosu" üzerinden performans önerileri listelenir.
\subsection{7. Açıklanabilir Yapay Zeka (SHAP Analizi)}
`07_shap_analizi.ipynb` notebook'unda `models/shap_values.pkl` içerisine kaydedilen verilerle beslenir. "Özellik Önemi", "Anomali Açıklama" ve "Model Karşılaştırma" sekmeleriyle sistemin neden alarm verdiğini açıklar (Örn: "Sensör X normal sınırın \%30 üzerinde").
\subsection{8. Anomali Detay ve PDF Raporlama}
Bir arıza segmenti seçildiğinde \texttt{reportlab.pdfgen} kütüphanesi kullanılarak, A4 boyutunda kurumsal bir "Anomali Raporu" PDF belgesi saniyeler içinde dinamik olarak çizilir (Canvas objeleriyle) ve kullanıcı bilgisayarına indirilir.
\subsection{9. Ablasyon (Ablation) Analizi Modülü}
\texttt{ablation\_page.py} sayfasıyla, özellik çıkarım denemeleri detaylandırılır:
\begin{itemize}
    \item \textbf{Tekil Özellik Etkisi:} Bir özellik çıkartıldığında AUC'nin nasıl etkilendiği.
    \item \textbf{Kümülatif Performans:} Eklenen her özelliğin getirdiği başarı artışı.
    \item \textbf{Grup Karşılaştırması:} 30-60-120 saniyelik pencerelerin grup başarım analizi.
    \item \textbf{Bağımlılık Matrisi:} Hangi özelliklerin birbirine bağımlı olduğunu gösteren Isı Haritası.
    \item \textbf{Önerilen Kritik Özellik Seti:} Toplam özelliklerin Pie Chart ile yüzdelik verimi gösterilir.
\end{itemize}

\section{Yazılım Mimari ve Mühendislik Pratikleri}
Proje, saf yapay zeka sınırlarının dışına çıkıp modern \textbf{DevOps ve Yazılım Mühendisliği} pratiklerini içerir:
\begin{itemize}
    \item \textbf{Bağımlılık Yönetimi (Dependencies):} `requirements.txt` ve `environment.yml` olmak üzere hem Pip hem de Conda ekosistemini destekleyen sanal ortam konfigürasyonu.
    \item \textbf{Birim Testleri (PyTest):} \texttt{tests/test\_data\_loader.py} ve \texttt{tests/test\_model\_loader.py} sayesinde, veri yüklendiğinde oluşabilecek şekilsel (shape) hatalar, boş argümanlar ve model yükleme hataları "Continuous Integration" mantığı ile test edilir.
    \item \textbf{TF Ortam Sabitleri:} \texttt{app.py} dosyasında işletim sistemi uyumsuzluklarını önlemek için \texttt{TF\_CPP\_MIN\_LOG\_LEVEL}, \texttt{CUDA\_VISIBLE\_DEVICES="-1"} gibi donanımsal sabitleyiciler eklenmiştir.
    \item \textbf{Otomasyon Betikleri:} \texttt{patch\_notebook.py} ve \texttt{revert\_notebook.py} gibi dosyalar kullanılarak projeyi kuran yeni mühendislerin Jupyter süreçlerini kod ile otomatik manipüle etmesi (A/B testing için) sağlanmıştır.
\end{itemize}

\section{Deneysel Sonuçlar ve Başarım}
Çapraz doğrulama (cross-validation) ve test seti sonuçlarına göre:
\begin{itemize}
    \item \textbf{Gözetimli Modellerde} çoklu katmanlı \textbf{MLP}, zaman alanı ve fiziksel sentetik metrikleri kullanarak mükemmele yakın bir doğruluğa ulaşmış (\%99.2 AUC, \%91.9 F1); sadece \%1.8 Yanlış Alarm Oranı (FAR) sergilemiştir. XGBoost ve Stacking Ensemble \%99.0 AUC ile yakın takipte kalmıştır.
    \item \textbf{Gözetimsiz (Etiketsiz) Analizde} bilinmeyen donanım anormalliklerini çözme yarışında derin öğrenme mimarisi \textbf{Autoencoder} (\%89.3 AUC) en iyi sonucu vermiştir.
    \item \textbf{Ablasyon Analizi Sonuçları:} On binlerce satırlık veri sütununun (time-domain, diff, lag) sadece en iyi \textbf{19 kritik özelliğe} indirgenmesinin sistem başarımının \%99'unu koruduğu matematiksel olarak kanıtlanmıştır. Uydu üzerindeki CPU/Bellek (SWaP-C) kısıtlamaları için muazzam bir tasarruf yaratılmıştır.
\end{itemize}

\section{Sonuç ve Kapanış}
ESA OPS-SAT telemetri verileriyle inşa edilen bu platform, sadece bir makine öğrenmesi modeli denemesi değil, modüler yapısı, otomatik raporlama mekanizması, canlı veri simulasyonu, API testleri ve açıklanabilirliğiyle (SHAP) doğrudan uydu yer kontrol istasyonlarına (Ground Control Station) entegre edilebilecek ticari/askeri seviyede bir FDIR (Arıza Tespit ve Giderme) ekosistemidir. 

\end{document}
